% PDF page 332
\source{gilbarg-trudinger:page-0332}

and the arguments of $D_{p_j}A^i$, $\delta_k A^i$ and $B$ are $x$, $u(x)$, $Du(x)$. Hence, writing

\[
\bar{a}^{ij}(x)=D_{p_j}A^i(x,u(x),Du(x))
\]
\[
f_k^i(x)=\delta_k A^i(x,u(x),Du(x))+\delta_k^iB(x,u(x),Du(x))
\]

where $[\delta_k^i]$ is the identity matrix, we see that the derivative $w=D_ku$ satisfies

\begin{equation}
\int_{\Omega} \bigl(\bar{a}^{ij}(x)D_jw+f_k^i(x)\bigr)D_i\zeta\,dx=0 \qquad \forall \zeta\in C_0^1(\Omega);
\tag{13.3}
\end{equation}

that is, $w$ is a generalized solution of the linear elliptic equation

\[
Lw=D_i(\bar{a}^{ij}D_jw)=-D_if_k^i.
\]

By replacing $\Omega$ if necessary by a strictly contained subdomain we can assume that
$L$ is strictly elliptic in $\Omega$ and that the coefficients $\bar{a}^{ij}$, $f_k^i$ are bounded, that is the
hypotheses of Theorems 8.22 and 8.24 are satisfied. Accordingly, choosing
$\lambda_K$, $\Lambda_K$, $\mu_K$ such that

\[
0<\lambda_K\leq \lambda(x,z,p),
\]
\begin{equation}
\Lambda_K\geq |D_{p_j}A^i(x,z,p)|,
\qquad
\mu_K\geq |\delta_jA^i(x,z,p)|+|B(x,z,p)|,
\tag{13.4}
\end{equation}

for all $x\in\Omega$, $|z|+|p|\leq K$, $i,j=1,\dots,n$, we obtain the following interior estimate.

\begin{theorem}\label{thm:gt13-holder-gradient-13-1}\dcref{Theorem 13.1}\group{ch13-holder-gradient}\source{gilbarg-trudinger:page-0332}
\textbf{Theorem 13.1.} \textit{Let $u\in C^2(\Omega)$ satisfy $Qu=0$ in $\Omega$ where $Q$ is elliptic in $\Omega$ and is
of divergence form (13.2) with $A\in C^1(\Omega\times\mathbb{R}\times\mathbb{R}^n)$, $B\in C^0(\Omega\times\mathbb{R}\times\mathbb{R}^n)$. Then for any
$\Omega'\Subset\Omega$ we have the estimate}
\end{theorem}

\begin{equation}
[Du]_{\alpha,\Omega'}\leq Cd^{-\alpha}
\tag{13.5}
\end{equation}

\textit{where}

\[
C=C(n,K,\Lambda_K/\lambda_K,\mu_K/\lambda_K,\operatorname{diam}\Omega),
\]
\[
K=|u|_{1,\Omega}=\sup_{\Omega}(|u|+|Du|),
\]
\[
d=\operatorname{dist}(\Omega',\partial\Omega)\text{ and }\alpha=\alpha(n,\Lambda_K/\lambda_K)>0.
\]

In order to extend Theorem 13.1 to a global H\"older estimate in $\Omega$ we assume
that $Q$ is elliptic in $\overline{\Omega}$ with $A\in C^1(\overline{\Omega}\times\mathbb{R}\times\mathbb{R}^n)$, $B\in C^0(\overline{\Omega}\times\mathbb{R}\times\mathbb{R}^n)$, that $\partial\Omega\in C^2$
and that $u=\varphi$ on $\partial\Omega$ where $\varphi\in C^2(\overline{\Omega})$. By replacing $u$ with $u-\varphi$ we can assume
without loss of generality that $u=0$ on $\partial\Omega$. Since $\partial\Omega\in C^2$, there exists for each
$x_0\in\partial\Omega$ a ball $B=B(x_0)$ and a one-to-one mapping $\psi$ from $B$ onto an open set

% PDF page 333
\source{gilbarg-trudinger:page-0333}

$D \subset \R^n$ such that

\[
\psi(B \cap \Omega) \subset \R^n_+ = \{x \in \R^n \mid x_n > 0\},
\qquad
\psi(B \cap \partial \Omega) \subset \partial \R^n_+,
\quad \text{and} \quad
\psi \in C^2(B), \ \psi^{-1} \in C^2(D).
\]

Writing $y=\psi(x)$, $v(y)=u \circ \psi^{-1}(y)$, $B^+=B\cap\Omega$, $D^+=\psi(B^+)$, we have $D_{y_k}v=0$ on $\partial D^+ \cap \partial \R^n_+$, $k=1,\ldots,n-1$, and that the equation $Qu=0$ in $B^+$ is equivalent to the equation

\begin{equation}\tag{13.6}
\bar Q v = D_{y_i}\bar A^i(x,u,Du) + \bar B(x,u,Du)=0, \qquad x=\psi^{-1}(y),
\end{equation}

in $D^+$ where the functions $\bar A$ and $\bar B$ are given by

\[
\bar A^i = \frac{\partial y_i}{\partial x_r} A^r, \qquad
\bar B = -\frac{\partial}{\partial y_i}\!\left(\frac{\partial y_i}{\partial x_r}\right) A^r + B.
\]

The derivatives $w=D_{y_k}v$, $k=1,\ldots,n$, will consequently be generalized solutions in $D^+$ of the linear elliptic equations

\[
Lw = D_i(\bar a^{ij}D_j w) = -D_i f_k^i,
\]

where

\[
\bar a^{ij}(y)=\frac{\partial y_i}{\partial x_r}\frac{\partial y_j}{\partial x_s}D_{p_s}A^r,
\]

\[
f_k^i(y)=\frac{\partial y_i}{\partial x_r}\frac{\partial x_s}{\partial y_k}
\left(\frac{\partial^2 y_l}{\partial x_j \partial x_s}D_{p_j}A^r D_{y_l}u+\delta_s A^r\right)+\delta_k^i B
\]
\[
\qquad + \left[\frac{\partial}{\partial y_k}\!\left(\frac{\partial y_i}{\partial x_r}\right)-\delta_k^i\frac{\partial}{\partial y_j}\!\left(\frac{\partial y_j}{\partial x_r}\right)\right]A^r,
\]

the arguments of $D_{p_r}A^r$, $\delta_s A^r$, $A^r$ and $B$ being $x=\psi^{-1}(y)$, $u(x)=v(y)$ and $Du(x)=Du(\psi^{-1}(y))$. By replacing $B(x_0)$ if necessary by a smaller concentric ball, we can then assume that the Jacobian matrix $[D\psi]=[\partial y_i/\partial x_j]$ is bounded from above and below in $B$ by positive constants, and consequently the operator $L$ is strictly elliptic in $D^+$ with bounded coefficients $\bar a^{ij}, f_k^i$. By applying Theorem 8.29, we thus have for any $D' \Subset D$

\begin{equation}\tag{13.7}
[D_{y_k}v]_{x;D'\cap D^+}\le C,\qquad k=1,\ldots,n-1,
\end{equation}

where $C=C(n,K,\Lambda_k/\lambda_k,\mu_k/\lambda_k,\Omega,d)$, $K=|u|_{1;\Omega}$, $d=\mathrm{dist}\,(D'\cap D^+,\partial D)$ and $\alpha=\alpha(n,\Lambda_k/\lambda_k,\Omega)>0$.

The remaining derivative $D_{y_n}v$ can be estimated as follows. Let $y_0\in D^+\cap D'$, $R\le d/3$, $B_{2R}=B_{2R}(y_0)$, $\eta\in C^1_0(B_{2R})$, and let $c$ be a constant such that $c=w(y_0)$ if $B_{2R}\subset D^+$, $c=0$ if $B_{2R}\cap \partial \R^n_+\ne \emptyset$. The function $\zeta=\eta^2(w-c)$ then belongs to $W^{1,2}_0(D^+)$ for $w=D_{y_k}v$, $k=1,\ldots,n-1$. By substitution into the integral identity


% PDF page 334
\source{gilbarg-trudinger:page-0334}

(13.3) with $\Omega = D^+$, we then have

\[
\int_{D^+} \eta^2 \tilde a^{ij} D_i w D_j w\,dy
\leq \int_{D^+} \Bigl\{|2\eta(w-c)\tilde a^{ij}D_i\eta D_j w|
+ |\eta^2 f^i D_i w| + |2\eta(w-c)f^iD_i\eta|\Bigr\}\,dy,
\]

so that by the Schwarz inequality (7.6) and the ellipticity of $L$, we obtain

\[
\int_{D^+} \eta^2|Dw|^2\,dy \leq C \int_{D^+} (\eta^2 + |D\eta|^2)(w-c)^2\,dy
\]

where $C = C(n,K,\Lambda_K/\lambda_K,\mu_K/\lambda_K,\Omega)$. Now let us further require $\eta$ to be such that
$0 \leq \eta \leq 1$, $\eta = 1$ in $B_R = B_R(y_0)$ and $|D\eta| \leq 2/R$. We thus obtain

\[
\int_{B_R} |Dw|^2\,dy \leq C R^{n-2}(R^2 + \sup_{B_{2R}}(w-c)^2)
\]

\[
\leq C R^{n-2+2\alpha}
\]

by (13.7), where $C$ and $\alpha$ depend on the same quantities as in (13.7). Therefore,
since $k=1,\dots,n-1$, we have

\begin{equation}
\int_{B_R} |D_i w|^2\,dy \leq C R^{n-2+2\alpha}
\tag{13.8}
\end{equation}

provided $j \neq n$. To proceed further we solve equation (13.6) for $D_{nn}v$, so that we
can write

\[
D_{nn}v = b^{ij}D_{ij}v + b,\qquad i=1,\dots,n,\quad j=1,\dots,n-1,
\]

for certain functions $b^{ij}$, $b$ bounded in terms of $Dw$, $K$, $\Lambda_K/\lambda_K$ and $\mu_K/\lambda_K$. Hence by
(13.8) we have

\[
\int_{B_R} |D_{ni}v|^2\,dy \leq C R^{n-2+2\alpha},\qquad i=1,\dots,n,
\]

so that using inequality (7.8) and Morrey's estimate, Theorem 7.19, we can
conclude that the estimate (13.7) is also valid for $k=n$. Returning to the domain $\Omega$
by means of the mapping $\psi^{-1}$, we thus have

\begin{equation}
\bigl[Du\bigr]_{x;B'\cap\Omega} \leq C
\tag{13.9}
\end{equation}

% PDF page 335
\source{gilbarg-trudinger:page-0335}


for any concentric ball $B' \subset B$, where $C = C(n, K, \Lambda_K/\lambda_K, \mu_K/\lambda_K, \Omega, B')$. Finally, by choosing a finite number of points $x_0 \in \partial\Omega$ and balls $B'$ covering $\partial\Omega$, we obtain the following global H\"older estimate from (13.5) and (13.9).

\begin{theorem}\label{thm:gt13-holder-gradient-13-2}\dcref{Theorem 13.2}\group{ch13-holder-gradient}\source{gilbarg-trudinger:page-0335}
\textbf{Theorem 13.2.} \emph{Let $u \in C^2(\overline{\Omega})$ satisfy $Qu = 0$ in $\Omega$ where $Q$ is elliptic in $\overline{\Omega}$ and is of divergence form (13.2) with $A \in C^1(\overline{\Omega}\times \mathbb{R}\times \mathbb{R}^n)$, $B \in C^0(\overline{\Omega}\times \mathbb{R}\times \mathbb{R}^n)$. Then if $\partial\Omega \in C^2$ and $u=\varphi$ on $\partial\Omega$, where $\varphi \in C^2(\overline{\Omega})$, we have the estimate}
\end{theorem}

\[
(13.10)\qquad [Du]_{x;\Omega} \leq C
\]

\emph{where}

\[
C = C(n, K, \Lambda_K/\lambda_K, \mu_K/\lambda_K, \Omega, \Phi),
\qquad K = |u|_{1;\Omega}, \qquad \Phi = |\varphi|_{2;\Omega}
\]

\[
\emph{and } \alpha = \alpha(n, \Lambda_K/\lambda_K, \Omega) > 0.
\]

An inspection of the proofs of Theorems 13.1 and 13.2 shows that the estimates (13.5) and (13.10) continue to hold if we only assume $u \in C^{0,1}(\overline{\Omega}) \cap W^{2,2}(\Omega)$, and $\varphi \in W^{2,q}(\Omega)$ for some $q>n$. In this generality we now must take $\Phi = \|\varphi\|_{W^{2,q}(\Omega)}$, and $\alpha$ will depend as well on $q$.

\section*{13.2. Equations in Two Variables}

If $u \in C^2(\Omega)$ satisfies the elliptic equation (13.1) in $\Omega \subset \mathbb{R}^2$, then the derivatives
$w_1 = D_1u$, $w_2 = D_2u$ are generalized solutions in $\Omega$ of the linear elliptic equations

\begin{align*}
L_1w_1 &= D_1\!\left(\frac{a^{11}}{a^{22}}D_1w_1 + \frac{2a^{12}}{a^{22}}D_2w_2\right) + D_{22}w_1 = -D_1\frac{b}{a^{22}}, \\
L_2w_2 &= D_{11}w_2 + D_2\!\left(\frac{2a^{12}}{a^{11}}D_1w_2 + \frac{a^{22}}{a^{11}}D_2w_2\right) = -D_2\frac{b}{a^{11}}.
\end{align*}

\[
(13.11)
\]

Consequently the methods for equations of divergence form also apply here. Accordingly, if $\lambda_K$, $\Lambda_K$ and $\mu_K$ satisfy

\[
0 < \lambda_K < \lambda(x,z,p),
\]
\[
(13.12)\qquad \Lambda_K \geq |a^{ij}(x,z,p)|,
\]
\[
\mu_K \geq |b(x,z,p)|,
\]

for all $x \in \Omega$, $|z|+|p| \leq K$, $i,j=1,2$, we have the following estimates.

\begin{theorem}\label{thm:gt13-holder-gradient-13-3}\dcref{Theorem 13.3}\group{ch13-holder-gradient}\source{gilbarg-trudinger:page-0335}
\textbf{Theorem 13.3.} \emph{Let $u \in C^2(\Omega)$ satisfy $Qu=0$ in $\Omega \subset \mathbb{R}^2$, where $Q$ is elliptic in $\Omega$ and the coefficients $a^{ij}$, $b \in C^0(\Omega\times \mathbb{R}\times \mathbb{R}^2)$. Then for any $\Omega' \subset \subset \Omega$, we have}
\end{theorem}

\[
(13.13)\qquad [Du]_{\alpha;\Omega'} \leq Cd^{-\alpha},
\]


% PDF page 336
\source{gilbarg-trudinger:page-0336}


where

\[
C = C(K,\Lambda_K/\lambda_K,\mu_K/\lambda_K,\operatorname{diam}\Omega),
\qquad K = |u|_{1;\Omega}, \qquad d = \operatorname{dist}(\Omega',\partial\Omega)
\quad \text{and} \quad \alpha = \alpha(\Lambda_K/\lambda_K) > 0.
\]

\begin{theorem}\label{thm:gt13-holder-gradient-13-4}\dcref{Theorem 13.4}\group{ch13-holder-gradient}\source{gilbarg-trudinger:page-0336}
\textbf{Theorem 13.4.} \emph{Let $u \in C^2(\overline{\Omega})$ satisfy $Qu = 0$ in $\Omega \subset \R^2$ where $Q$ is elliptic in $\overline{\Omega}$ and the coefficients $a^{ij}, b \in C^0(\overline{\Omega} \times \R \times \R^2)$. Then if $\partial\Omega \in C^2$, $Q \in C^2(\overline{\Omega})$ and $u = \varphi$ on $\partial\Omega$, we have the estimate}
\end{theorem}

\[
(13.14)\qquad [Du]_{x;\Omega} \leq C
\]

\emph{where}

\[
C = C(K,\Lambda_K/\lambda_K,\mu_K/\lambda_K,\Omega,\Phi),
\qquad K = |u|_{1;\Omega}, \qquad \Phi = |\varphi|_{2;\Omega}
\]

\[
\emph{and} \quad \alpha = \alpha(\Lambda_K/\lambda_K,\Omega) > 0.
\]

Note that we have given an alternative derivation of Theorem 13.3 in Chapter
12 using the method of quasiconformal mappings. We note also that in the case of
two variables the proof of the H\"older estimate for linear divergence form equations
is simpler than for more variables (see Problem 8.5). The remark following Theorem
13.2 of course applies as well to Theorems 13.3 and 13.4.

\section*{13.3. Equations of General Form; the Interior Estimate}

We shall treat elliptic equations of the general form (13.1) by showing that certain
combinations of the derivatives of solutions are generalized subsolutions of linear
elliptic equations of divergence form. The desired H\"older estimates are then
obtained through application of the weak Harnack inequalities, Theorems 8.18
and 8.26.

Let us suppose that the coefficients $a^{ij}$ and $b$ of $Q$, are respectively in
$C^1(\Omega \times \R \times \R^n)$ and $C^0(\Omega \times \R \times \R^n)$. Let $Qu = 0$ in
$\Omega$ and assume initially that $u \in C^3(\Omega)$. By differentiation with respect to
$x_k$, $k = 1,\dots,n$, we obtain

\[
(13.15)\qquad a^{ij}(x,u,Du)D_{kij}u + D_{p_l}a^{ij}(x,u,Du)D_{lk}u\,D_{ij}u
\]
\[
\qquad\qquad + \delta_k a^{ij}(x,u,Du)D_{ij}u + D_k b(x,u,Du) = 0
\]

where $\delta_k$ is the differential operator defined by

\[
\delta_k g(x,z,p) = D_{x_k}g(x,z,p) + p_k D_z g(x,z,p).
\]

The equation (13.15) can be written in the following divergence form

\[
(13.16)\qquad D_i(a^{ij}D_{kj}u) + (D_{p_l}a^{ij} - D_{p_j}a^{il})D_{lk}uD_{ij}u
\]
\[
\qquad\qquad + \delta_k a^{ij}D_{ij}u - \delta_i a^{ij}D_{kj}u + D_k b = 0.
\]


% PDF page 337
\source{gilbarg-trudinger:page-0337}


Let us now write $v=|Du|^2$, multiply equation (13.16) by $D_ku$ and sum the resulting equations from $k=1$ to $n$. We then obtain

\begin{equation}
\begin{aligned}
&-a^{ij}D_{ki}uD_{kj}u+\frac{1}{2}D_i(a^{ij}D_jv)+\frac{1}{2}(D_p a^{ij}-D_p a^{il})D_i uD_jv \\
&\qquad +D_ku\,\delta_k a^{ij}D_{ij}u-\frac{1}{2}\delta_k a^{ij}D_jv+D_k(bD_ku)-b\,\Delta u=0.
\end{aligned}
\tag{13.17}
\end{equation}

For $\gamma\in\R$ and $r=1,\dots,n$, we next define functions

\begin{equation}
w=w_r=\gamma D_ru+v,
\tag{13.18}
\end{equation}

and by combining (13.16) and (13.17) obtain the equation,

\begin{equation}
\begin{aligned}
&-2a^{ij}D_{ki}uD_{kj}u+D_i\!\bigl(a^{ij}D_jw+(2D_iu+\gamma\delta^{ir})b\bigr) \\
&\quad +(D_{p_i}a^{ij}-D_{p_j}a^{il})D_{ij}uD_lw+\bigl[(2D_ku+\gamma\delta^{kr})\delta_k a^{ij}-2b\delta^{ij}\bigr]D_{ij}u-\delta_i a^{ij}D_jw \\
&\qquad =0.
\end{aligned}
\tag{13.19}
\end{equation}

Setting

\[
\begin{aligned}
\bar a^{ij}(x)&=a^{ij}(x,u(x),Du(x)),\\
a_i^{\,j}(x)&=(D_{p_i}a^{ij}-D_{p_j}a^{il})(x,u(x),Du(x)),\\
f_i(x)&=(2D_iu(x)+\gamma\delta^{ir})b(x,u(x),Du(x)),\\
b^{ij}(x)&=\bigl[(2D_ku(x)+\gamma\delta^{kr})\delta_k a^{ij}-2\delta^{ij}b\bigr](x,u(x),Du(x)),\\
c^j(x)&=\delta_i a^{ij}(x,u(x),Du(x)),
\end{aligned}
\]

we write the equation (13.19) in its integral form

\begin{equation}
\int_\Omega \Bigl\{(\bar a^{ij}D_jw+f^i)D_i\zeta+\bigl(2\bar a^{ij}D_{ki}uD_{kj}u-a_i^{\,j}D_{ij}uD_lw+c^iD_iw-b^{ij}D_{ij}u\bigr)\zeta\Bigr\}\,dx
\tag{13.20}
\end{equation}

\[
=0
\]

for all $\zeta\in C_0^1(\Omega)$. We claim now that the integral identity (13.20) continues to be valid if we only assume that $u\in C^2(\Omega)$. To see this, we let $\{u_m\}\subset C^3(\Omega)$ approach $u$ in the sense of $C^2(\Omega)$, that is $\{D^\beta u_m\}$ converges uniformly to $D^\beta u$ on compact subsets of $\Omega$ for all $|\beta|\leq 2$. Since $Qu=0$ in $\Omega$, we have $Qu_m\to 0$ uniformly on compact subsets of $\Omega$ and hence letting $m\to\infty$ in the integral identity for $u_m$ corresponding to (13.20), we obtain (13.20). We note that a similar approximation argument shows that in fact we need only assume $u\in C^{0,1}(\Omega)\cap W^{2,2}(\Omega)$.

To proceed further, we need to remove the terms involving $D_{ij}u$ from (13.20).
Since $Q$ is elliptic we have

\[
\bar a^{ij}D_{ki}uD_{kj}u\geq \lambda|D^2u|^2\geq 0
\]

where $\lambda=\lambda(x,u(x),Du(x))$. Using the Schwarz inequality, we then have from (13.20)


% PDF page 338
\source{gilbarg-trudinger:page-0338}


\[
(13.21)\qquad \int_{\Omega} (\tilde a^{ij}D_j w + \tilde f^i)D_i\zeta \, dx
\]
\[
\leq \int_{\Omega} \left\{\left(\lambda + \frac{1}{\lambda}\sum |a^{ij}|^2\right)|Dw|^2 + \frac{1}{\lambda}\sum \left(|c^i|^2 + |b^{ij}|^2\right)\right\}\zeta \, dx
\]

for all non-negative $\zeta \in C_0^1(\Omega)$. The reduction to the linear theory is finally achieved
by replacing $\zeta$ in (13.21) by $\zeta e^{2\chi w}$ where $\chi = \sup_{\Omega} \left(1 + \lambda^{-2}\sum |a^{ij}|^2\right)$. Setting

\[
\tilde a^{ij}(x) = e^{2\chi w(x)} a^{ij}(x)
\]
\[
\tilde f^i(x) = e^{2\chi w(x)} f^i(x)
\]
\[
\tilde g(x) = \frac{1}{\lambda} e^{2\chi w(x)} \left(\chi \sum |f^i|^2 + \sum |c^i|^2 + \sum |b^{ij}|^2\right)
\]

we thus obtain

\[
(13.22)\qquad \int_{\Omega} \{(\tilde a^{ij}D_i w + \tilde f^j)D_j\zeta - \tilde g\zeta\}\, dx \leq 0
\]

for all non-negative $\zeta \in C_0^1(\Omega)$; that is, the function $w$ satisfies the inequality

\[
(13.23)\qquad Lw = D_i(\tilde a^{ij}D_j w) \geq -(\tilde g + D_i\tilde f^i)
\]

in the generalized sense. By replacing $\Omega$ if necessary by a strictly contained sub-
domain we can assume that the operator $L$ is strictly elliptic in $\Omega$ and that the
coefficients $\tilde a^{ij}$, $\tilde f^i$ and $\tilde g$ are bounded.

    Let us take $\gamma > 0$ and write $w_r^{\pm} = \pm \gamma D_r u + v$. We now show that the validity of
the inequality (13.23) for all sufficiently large $\gamma \in \R$ and $r = 1,\ldots,n$ is sufficient
for the derivation of H\"older estimates for the derivatives $D_r u$, $r = 1,\ldots,n$. For, by
choosing $\gamma$ sufficiently large, we can ensure that the functions $w_r^{\pm}$ behave in a
certain sense the same as $\pm D_r u$. It turns out that a convenient choice of $\gamma$ is
$\gamma = 10nM$ where $M = \sup |Du|$. For with this choice if $\mathcal{G}$ is an arbitrary subset of $\Omega$
and $r$ is chosen so that

\[
\osc_{\mathcal{G}} D_r u \geq \osc_{\mathcal{G}} D_i u, \quad i = 1,\ldots,n,
\]

it is readily seen that

\[
(13.24)\qquad 8nM\, \osc_{\mathcal{G}} D_r u \leq \osc_{\mathcal{G}} w_r^{\pm} \leq 12nM\, \osc_{\mathcal{G}} D_r u.
\]

Furthermore, writing $w^{\pm} = w_r^{\pm}$ for brevity and setting $W^{\pm} = \sup_{\mathcal{G}} w^{\pm}$, we have


% PDF page 339
\source{gilbarg-trudinger:page-0339}

\setcounter{section}{13}
\setcounter{subsection}{2}
\noindent 13.3. Equations of General Form; the Interior Estimate \hfill 327

\begin{equation}\tag{13.25}
\inf_{\zeta\in\mathcal{G}}\sum_{\pm}\left(W^{\pm}-w^{\pm}\right)\geq 10n\,M\left(\sup_{\mathcal{G}}D_ru-\inf_{\mathcal{G}}D_ru\right)+2\inf_{\mathcal{G}}v-2\sup_{\mathcal{G}}v
\end{equation}
\[
\geq 6n\,M\,\osc_{\mathcal{G}}D_ru
\]
\[
\geq \frac{1}{2}\,\osc_{\mathcal{G}}w^{\pm}.
\]

We are now in a convenient position to apply the weak Harnack inequality, Theorem 8.18. Let us take $\mathcal{G}=B_{4R}(y)\subset \Omega$. The functions $u=W^{\pm}-w^{\pm}$ will be non-negative supersolutions in $\mathcal{G}$ of the equation
\[
Lu=-(\sg+D_i f^i).
\]

Accordingly, choosing $\lambda_k$ and $\mu_k$ such that
\[
0<\lambda_k<\lambda(x,z,p),
\]
\begin{equation}\tag{13.26}
\mu_k\geq |a^{ij}(x,z,p)|+\left|D_{p_k}a^{ij}(x,z,p)\right|+\left|D_z a^{ij}(x,z,p)\right|
\end{equation}
\[
+\left|D_{x_k}a^{ij}(x,z,p)\right|+|b(x,z,p)|,
\]
for all $x\in \Omega$, $|z|+|p|\leq K$, $i,j,k=1,\ldots,n$, and setting $p=1$ in Theorem 8.18, we obtain the estimates
\begin{equation}\tag{13.27}
R^{-n}\int_{B_{2R}}(W^{\pm}-w^{\pm})\,dx\leq C\bigl(W^{\pm}-\sup_{B_R}w^{\pm}+\sigma(R)\bigr),
\end{equation}
where $C=C(n,K,\mu_k/\lambda_k)$, $K=|u|_{1;\Omega}$ and $\sigma(R)=R+R^2$. Using (13.25), we see that the inequality
\[
\frac{1}{\omega_n(2R)^n}\int_{B_{2R}}(W^{\pm}-w^{\pm})\,dx\geq \frac{1}{4}\,\osc_{B_{4R}}w^{\pm},
\]
holds for either $w^+$ or $w^-$. Let us suppose it is true for $w^+$. Then from (13.27) we have
\[
\osc_{B_{4R}}w^+\leq C\bigl(W^+-\sup_{B_R}w^++\sigma(R)\bigr)
\]
\[
\leq C\bigl(\osc_{B_{4R}}w^+-\osc_{B_R}w^++\sigma(R)\bigr),
\]
and consequently, writing $\omega(R)=\osc_{B_R}w^+$, we have
\[
\omega(R)\leq \gamma\omega(4R)+\sigma(R),
\]
where $\gamma=1-C^{-1},\ C=C(n,K,\mu_k/\lambda_k)$.

We need now the following extension of Lemma 8.23.


% PDF page 340
\source{gilbarg-trudinger:page-0340}

\begin{lemma}\label{lem:gt13-holder-gradient-13-5}\dcref{Lemma 13.5}\group{ch13-holder-gradient}\source{gilbarg-trudinger:page-0340}
Lemma 13.5. \emph{Let $\omega_1,\dots,\omega_N$, $\bar{\omega}_1,\dots,\bar{\omega}_N\ge 0$ be non-decreasing functions on an interval $(0,R_0)$ such that for each $R\le R_0$ a function $\bar{\omega}_r$ can be found satisfying the inequalities}

\[
\bar{\omega}_r(R)\ge \delta_0\omega_i(R), \qquad i=1,\dots,N, \quad \text{and}
\]
\begin{equation}\tag{13.28}
\bar{\omega}_r(\delta R)\le \gamma \bar{\omega}_r(R)+\sigma(R)
\end{equation}

\emph{where $\sigma$ is also non-decreasing, $\delta_0>0$ and $0<\gamma$, $\delta<1$. Then for any $\mu\in(0,1)$ and $R\le R_0$, we have}

\begin{equation}\tag{13.29}
\omega_i(R)\le C\left\{\left(\frac{R}{R_0}\right)^\alpha \omega_0+\sigma(R^\mu R_0^{1-\mu})\right\}
\end{equation}

\emph{where $C=C(N,\delta_0,\delta,\gamma)$ and $\alpha=\alpha(N,\delta_0,\delta,\gamma,\mu)$ are positive constants and $\omega_0=\max_{i=1,\dots,N}\bar{\omega}_i(R_0)$.}
\end{lemma}

The proof of Lemma 13.5 follows by a simple modification of Lemma 8.23 and is therefore omitted. If we now let $B_0=B_{R_0}(y)$ be any ball contained in $\Omega$, it then follows by Lemma 13.5 with $N=2n$, $\delta_0=8nM$, $\delta=\frac14$, and $\mu=\frac12$, that for any $R\le R_0$, $i=1,\dots,n$

\begin{equation}\tag{13.30}
\osc_{B_R(y)} D_i u \le CR^\alpha
\end{equation}

where $C=C(n,K,\mu_K/\lambda_K,R_0)$, $\alpha=\alpha(n,K,\mu_K/\lambda_K)$. Consequently we obtain the desired interior estimate asserted in the following theorem.

\begin{theorem}\label{thm:gt13-holder-gradient-13-6}\dcref{Theorem 13.6}\group{ch13-holder-gradient}\source{gilbarg-trudinger:page-0340}
Theorem 13.6. \emph{Let $u\in C^2(\Omega)$ satisfy $Qu=0$ in $\Omega$ where $Q$ is elliptic in $\Omega$ and the coefficients $a^{ij}\in C^1(\Omega\times\mathbb{R}\times\mathbb{R}^n)$, $b\in C^0(\Omega\times\mathbb{R}\times\mathbb{R}^n)$. Then for any $\Omega'\Subset\Omega$ we have the estimate}

\begin{equation}\tag{13.31}
[Du]_{x,\Omega'}\le Cd^{-\alpha}
\end{equation}

\emph{where $C=C(n,K,\mu_K/\lambda_K,\operatorname{diam}\Omega)$, $K=|u|_{1;\Omega}$, $d=\dist(\Omega',\partial\Omega)$ and $\alpha=\alpha(n,K,\mu_K/\lambda_K)$.}
\end{theorem}

\vspace{1em}

\noindent 13.4. Equations of General Form; the Boundary Estimate

Let us suppose now that the operator $Q$ is elliptic in $\overline{\Omega}$ and that the coefficients $a^{ij}$ and $b$ are respectively in $C^1(\overline{\Omega}\times\mathbb{R}\times\mathbb{R}^n)$ and $C^0(\overline{\Omega}\times\mathbb{R}\times\mathbb{R}^n)$. Let $\partial\Omega\in C^2$ and assume $u=\varphi$ on $\partial\Omega$ where $\varphi\in C^2(\overline{\Omega})$. By replacing $u$ by $u-\varphi$, we can assume without loss of generality that $u=0$ on $\partial\Omega$. Furthermore, if $\mathcal{U}$ is an open subset of $\Omega$ and $x\mapsto y=\psi(x)$ defines a $C^2(\mathcal{U})$ coordinate transformation, we have for $x\in\mathcal{U}$

% PDF page 341
\source{gilbarg-trudinger:page-0341}

\setcounter{section}{13}
\setcounter{subsection}{3}
\noindent 13.4. Equations of General Form; the Boundary Estimate \hfill 329

\begin{equation}\tag{13.32}
Qu=a^{kl}(x,u,Du)\,\frac{\partial y_i}{\partial x_k}\,\frac{\partial y_j}{\partial x_l}\,D_{y_i y_j}u
+a^{ij}(x,u,Du)\,\frac{\partial^2 y_k}{\partial x_i\,\partial x_j}\,D_{y_k}u
+b(x,u,Du).
\end{equation}

Consequently the equation $Qu=0$ in $\Omega$ is transformed into an equation $\bar{Q}v=0$
in $\psi(\Omega)$ where $v=u\circ\psi^{-1}$ and $\bar{Q}$ satisfies the same hypotheses as $Q$. It therefore
suffices to consider equation (13.1) in the neighborhood of a flat boundary
portion. Accordingly, let $D$ be an open set in $\mathbf{R}^n$ such that $D^+=D\cap\Omega\subset\mathbf{R}^n_+=
\{x\in\mathbf{R}^n|x_n>0\}$ and $D\cap\partial\Omega=\partial\mathbf{R}^n_+$. We define functions $v'$ and $w$ by

\begin{equation}\tag{13.33}
v'=\sum_{i=1}^{n-1}|D_i u|^2,\qquad w=w_r=\gamma D_r u+v',\qquad r=1,\ldots,n-1,\ \gamma\in\mathbf{R}.
\end{equation}

It is evident that $w=0$ on $D\cap\partial\Omega$ and that $w$ satisfies equation (13.20) provided
the summation over the index $k$ is taken from $1$ to $n-1$. Under this restriction we
then have

\[
\bar{a}^{ij}D_{ki}u\,D_{kj}u\geq \lambda \sum_{j\neq n}|D_{ij}u|^2.
\]

The missing derivative $D_{nn}u$ is estimated by writing equation (13.1) in the form

\begin{equation}\tag{13.34}
D_{nn}u=-\frac{1}{a^{nn}}\left(\sum_{(i,j)\neq(n,n)}a^{ij}D_{ij}u+b\right).
\end{equation}

Inserting (13.34) into (13.20) and proceeding as in the preceding section, we arrive
at the integral inequality (13.22) with $\Omega$ replaced by $D^+$ and $\chi$ and $\sg$ replaced
respectively by

\[
\chi=\sup_{\Omega}\left(1+\lambda^{-2}\sum|a^{ij}|^2\right)\left(1+\lambda^{-2}\sum|a^{ij}|^2\right)
\]

and

\[
\sg=\lambda^{-1}e^{2xw}\left(\chi\sum|f^i|^2+\sum|c^i|^2+b^2\left(1+\lambda^{-2}\sum|a^{ij}|^2\right)
+\sum|b^{ij}|^2\left(1+\lambda^{-2}\sum|a^{ij}|^2\right)\right).
\]

By considering balls centered at $y\in D\cap\partial\Omega$, applying the boundary weak Harnack
inequality (Theorem 8.26) instead of the interior Harnack inequality (Theorem
8.18) and following the proof of the preceding section, we obtain a boundary
Hölder estimate for the tangential derivatives of $u$. That is, for any ball $B_0=
B_{R_0}(y)\subset D$ with center $y\in D\cap\partial\Omega$, we have for any $R\leq R_0$

\begin{equation}\tag{13.35}
\osc_{D^+\cap B_R(y)} D_i u\leq CR^\alpha,\qquad i=1,\ldots,n-1,
\end{equation}

where $C=C(n,K,\mu_K/\lambda_K,R_0)$, $K=|u|_{1;\Omega}$ and $\alpha=\alpha(n,K,\mu_K/\lambda_K)>0$.


% PDF page 342
\source{gilbarg-trudinger:page-0342}

The passage from (13.35) to an estimate for the remaining derivative $D_nu$ is similar to the corresponding step for divergence form equations in Section 13.1. Let $D'\subset\subset D$, $d=\operatorname{dist}(D'\cap D^+,\,\partial D)$, $y\in D^+\cap D'$, $R\le d/3$, $B_{2R}=B_{2R}(y)$, $\eta\in C_0^1(B_{2R})$, and let $c$ be a constant such that $c=\inf w$ if $B_{2R}\subset D^+$, $c=0$ if $B_{2R}\cap\partial\Omega\neq\phi$. The function $\zeta=\eta^2(w-c)^+=\eta^2\sup_{B_{2R}}(w-c,0)$ is non-negative and belongs to $W_0^{1,2}(D^+)$ for $w$ given by (13.33). By substitution into the integral inequality (13.22) with $\Omega=D^+$, we then have

\[
\int_{w\ge c}\eta^2|Dw|^2\,dx\le C\int_{D^+}(\eta^2+|D\eta|^2)(w-c)^2\,dx
\]

where $C=C(n,K,\mu_K/\lambda_K)$. Now let us further require $\eta$ to be such that $0\le\eta\le1$, $\eta=1$ in $B_R=B_R(y)$ and $|D\eta|\le 2/R$. We then obtain

\begin{equation}
\int_{B_R^+}|Dw|^2\,dx\le CR^{n-2}\bigl(R^2+\sup_{B_{2R}}(w-c)^2\bigr),\qquad B_R^+=\{x\in B_R\mid w(x)\ge c\},
\tag{13.36}
\end{equation}

\[
\le CR^{n-2}\bigl(R^2+(\osc_{B_{2R}}w)^2\bigr)
\]

\[
\le CR^{n-2+2\alpha}\qquad \text{by (13.30) and (13.35),}
\]

where $C=C(n,K,\mu_K/\lambda_K,d)$ and $\alpha=\alpha(n,K,\mu_K/\lambda_K)$. By taking $\gamma=0$ and $\gamma=1$ in (13.36), we then have for $B_{2R}\subset D^+$ (in which case $B_R^+=B_R=D^+\cap B_R$)

\begin{equation}
\int_{D^+\cap B_R}|DD_ru|^2\,dx\le 2\int_{D^+\cap B_R}\bigl(|Dv'|^2+|Dw|^2\bigr)\,dx,
\tag{13.37}
\end{equation}

\[
\le CR^{n-2+2\alpha},\qquad r=1,\ldots,n-1.
\]

If $B_{2R}\cap\partial\Omega\neq\phi$, we require $\gamma=0,\pm1$ in (13.36). The estimate (13.37) then follows again since at each point of $D'\cap B_R$ at least one of the functions $w^{\pm}=\pm D_ru+v'$ is non-negative. Therefore we have the estimate

\begin{equation}
\int_{D^+\cap B_R}|D_{ij}u|^2\,dx\le CR^{n-2+2\alpha}
\tag{13.38}
\end{equation}

for any $y\in D'\cap D^+$, $R\le d/3$, provided $j\neq n$. By virtue of (13.34), the estimate (13.38) is also valid for $i=j=n$. Hence, by Theorem 7.19, we have

\begin{equation}
[Du]_{x;D'\cap D^+}\le C
\tag{13.39}
\end{equation}

where $C=C(n,K,\mu_K/\lambda_K,d)$ and $\alpha=\alpha(n,K,\mu_K/\lambda_K)>0$. Finally, by returning to


% PDF page 343
\source{gilbarg-trudinger:page-0343}

the original domain $\Omega$ and boundary values $\varphi$ we obtain the fundamental global
H\"older estimate of Ladyzhenskaya and Ural'tseva.

\begin{theorem}\label{thm:gt13-holder-gradient-13-7}\dcref{Theorem 13.7}\group{ch13-holder-gradient}\source{gilbarg-trudinger:page-0343}
\textbf{Theorem 13.7.} \textit{Let $u \in C^2(\bar{\Omega})$ satisfy $Qu=0$ in $\Omega$ where $Q$ is elliptic in $\bar{\Omega}$ and the
coefficients $a^{ij} \in C^1(\bar{\Omega}\times \mathbb{R}\times \mathbb{R}^n)$, $b \in C^0(\bar{\Omega}\times \mathbb{R}\times \mathbb{R}^n)$. Then if $\partial \Omega \in C^2$, $\varphi \in C^2(\bar{\Omega})$
with $u=\varphi$ on $\partial \Omega$, we have the estimate}
\end{theorem}

\begin{equation}
(13.40)\qquad [Du]_{x;\Omega}\leq C
\end{equation}

where

\[
C=C(n,K,\mu_K/\lambda_K,\Omega,\Phi),
\]
\[
K=|u|_{1;\Omega},\ \Phi=|\varphi|_{2;\Omega}\qquad \text{and}\qquad \alpha=\alpha(n,K,\mu_K/\lambda_K,\Omega)>0.
\]

An inspection of the proofs of Theorems 13.6 and 13.7 shows that the estimates
(13.31) and (13.40) continue to hold if we only assume $u \in C^0(\bar{\Omega}) \cap W^{2,2}(\Omega)$ and
$\varphi \in W^{2,q}(\Omega)$ for some $q>n$. In this generality we now must take $\Phi=|\varphi|_{W^{2,q}(\Omega)}$ and
the constants $C$ and $\alpha$ will depend on $q$.

Furthermore, it is evident from the development of this chapter that the global
and interior H\"older estimates can be realized as special cases of a partially interior
estimate. Namely, let us suppose that the operator $Q$ satisfies the hypotheses of
Theorem 13.7 and let $T$ be a $C^2$ boundary portion of $\partial \Omega$ such that $u,\varphi \in C^2(\Omega\cup T)$
with $u=\varphi$ on $T$. Then, if $Qu=0$ in $\Omega$, we have for any $\Omega' \Subset \Omega\cup T$ the estimate

\begin{equation}
(13.41)\qquad [Du]_{x;\Omega'}\leq C,
\end{equation}

where

\[
C=C(n,K,\mu_K/\lambda_K,T,\Phi,d),
\]
\[
d=\operatorname{dist}(\Omega',\,\partial \Omega - T)\qquad \text{and}\qquad \alpha=\alpha(n,K,\mu_K/\lambda_K,T)>0.
\]

\section*{13.5. Application to the Dirichlet Problem}

By combining Theorems 11.4, 13.2, 13.4, 13.7, we obtain the following fundamental
existence theorem.

\begin{theorem}\label{thm:gt13-holder-gradient-13-8}\dcref{Theorem 13.8}\group{ch13-holder-gradient}\source{gilbarg-trudinger:page-0344}
\textbf{Theorem 13.8.} \textit{Let $\Omega$ be a bounded domain in $\mathbb{R}^n$ and suppose that the operator $Q$
is elliptic in $\bar{\Omega}$ with coefficients $a^{ij} \in C^1(\bar{\Omega}\times \mathbb{R}\times \mathbb{R}^n)$, $b \in C^\alpha(\bar{\Omega}\times \mathbb{R}\times \mathbb{R}^n)$, $0<\alpha<1$.
Let $\partial \Omega \in C^{2,\alpha}$ and $\varphi \in C^{2,\alpha}(\bar{\Omega})$. Then, if there exists a constant $M$, independent of $u$
and $\sigma$, such that every $C^{2,\alpha}(\bar{\Omega})$ solution of the Dirichlet problems,}
\end{theorem}

\begin{equation}
(13.42)\qquad Q_\sigma u=a^{ij}(x,u,Du) D_{ij}u+\sigma b(x,u,Du)=0 \text{ in } \Omega,
\end{equation}
\[
u=\sigma\varphi \text{ on } \partial \Omega,\qquad 0\leq \sigma\leq 1,
\]


% PDF page 344
\source{gilbarg-trudinger:page-0344}

satisfies

\[
(13.43)\qquad \|u\|_{C^1(\bar{\Omega})}=\sup_{\Omega}|u|+\sup_{\Omega}|Du|<M,
\]

\emph{it follows that the Dirichlet problem, $Qu=0$ in $\Omega$, $u=\varphi$ on $\partial\Omega$ is solvable in $C^2(\bar{\Omega})$.}
\emph{If $Q$ is of divergence form or if $n=2$ we need only assume $a^{ij}\in C^1(\bar{\Omega}\times \mathbb{R}\times \mathbb{R}^n)$.}

The solvability of the Dirichlet problem is reduced by Theorem 13.8 to the
apriori estimation in the space $C^1(\bar{\Omega})$ of a related family of problems. Provided the
hypotheses of Theorem 13.8 hold, we therefore need only carry out the first three
steps of the existence procedure described in Chapter 11, Section 3. Furthermore
by invoking the more general Theorem 11.8 instead of Theorem 11.4, we see that
the family of problems (13.42) can be replaced in the statement of Theorem 13.8 by
any family of the form

\[
(13.44)\qquad Q_\sigma u=a^{ij}(x,u,Du;\sigma)D_{ij}u+b(x,u,Du;\sigma)=0\ \text{in }\Omega,
\]
\[
u=\sigma\varphi\ \text{on }\partial\Omega,\qquad 0\le \sigma\le 1
\]

where

\begin{enumerate}
\item[(i)] $Q_1=Q$, $b(x,z,p;0)=0$.
\item[(ii)] the operators $Q_\sigma$ are elliptic in $\bar{\Omega}$ for all $\sigma\in[0,1]$,
\item[(iii)] the coefficients $a^{ij}$, $b$ are sufficiently smooth; for example, $a^{ij}$, $b\in$
$C^1(\bar{\Omega}\times \mathbb{R}\times \mathbb{R}^n\times[0,1])$.
\end{enumerate}

\textbf{Notes}

The Hölder estimates in Theorems 13.1, 13.2, 13.6 and 13.7 are due to Ladyzhen-
skaya and Ural'tseva, [LU 2, 3, 4]. Our derivation of Theorems 13.6 and 13.7,
adapted from [TR 1], differs somewhat from theirs, although we retain their key
idea of a reduction to a divergence structure inequality for the functions $w$ in
(13.33). Note that we may use the weak Harnack inequality, Theorem 9.22, in place
of divergence structure results in the derivation of Theorem 13.6 (see [TR 13]).
Divergence structure theory can also be avoided in Theorem 13.7 by means of
Krylov's boundary gradient Hölder estimate, Theorem 9.31; (see Problem 13.1).

\textbf{Problem 13.1.} Using the interpolation inequality, Lemma 6.32, show that the
interior estimates (13.5), (13.13), (13.31) can be cast more explicitly in the form

\[
(13.45)\qquad [Du]_{\alpha;\Omega'}\leqq C\bigl(d^{1-\alpha}\sup_{\Omega''}|u|+1\bigr)
\]

for any $\Omega'\Subset\Omega''\Subset\Omega$ where $C$ and $\alpha$ depend on the same quantities as before
and $d=\operatorname{dist}(\Omega',\partial\Omega'')$. By combining (13.45) with the boundary estimate (9.68)
(analogously to the proof of Theorem 8.29) deduce the global estimates of (13.10),
(13.14), (13.40) for solutions $u\in C^{0,1}(\bar{\Omega})\cap C^2(\Omega)$. We remark that these global
estimates may also be derived directly from the present forms of (13.5), (13.13),
(13.31); (see [TR 14]).

% PDF page 345
\source{gilbarg-trudinger:page-0345}

\setcounter{chapter}{13}
\noindent Chapter 14

\bigskip

\noindent\textbf{Boundary Gradient Estimates}

\bigskip

An examination of the proof of Theorem 11.5 shows that for elliptic operators of
the forms (11.7) or (11.8) the solvability of the classical Dirichlet problem with
smooth data depends only upon the fulfillment of Step II of the existence pro-
cedure, that is, upon the existence of a boundary gradient estimate. In this chapter
we provide a variety of hypotheses for the general equation,

\begin{equation}
Qu=a^{ij}(x, u, Du)D_{ij}u+b(x, u, Du)=0
\tag{14.1}
\end{equation}

in $\Omega\subset\mathbb{R}^n$, that guarantee a boundary gradient estimate for solutions. These
hypotheses are combinations of structural conditions on $Q$ and geometric con-
ditions on the domain $\Omega$. It will be seen that the gradient bound aspect of the theory
of quasilinear elliptic equations is not as profound as other aspects such as the
Hölder estimates of Chapters 6 and 13. The boundary gradient estimates are tied
through the classical maximum principle to judicious and generally natural choices
of barrier functions. Nevertheless these estimates are of considerable importance
since they seem to be the principal factor in determining the solvability character
of the Dirichlet problem. This will be evidenced by the non-existence results at the
end of the chapter.

\medskip

    A description of the barrier method to be employed below is appropriate at this
juncture. This is a modification of ideas already met in Chapters 2 and 6. Let $Q$
be an elliptic operator of the form,

\begin{equation}
Qu=a^{ij}(x, Du)D_{ij}u+b(x, u, Du),
\tag{14.2}
\end{equation}

where $b(x, z, p)$ is non-increasing in $z$, and suppose that $u\in C^2(\Omega)\cap C^0(\overline{\Omega})$
satisfies $Qu=0$ in $\Omega$. Suppose that, in some neighborhood $\mathcal{N}=\mathcal{N}_{x_0}$ of a point
$x_0\in\partial\Omega$, there exist two functions $w^\pm=w^\pm_{x_0}\in C^2(\mathcal{N}\cap\Omega)\cap C^1(\mathcal{N}\cap\overline{\Omega})$ such that

\begin{itemize}
\item[(i)] $\pm Qw^\pm<0$ in $\mathcal{N}\cap\Omega$
\item[(ii)] $w^\pm(x_0)=u(x_0)$
\item[(iii)] $w^-(x)\leq u(x)\leq w^+(x), \qquad x\in\partial(\mathcal{N}\cap\Omega)$.
\end{itemize}

It then follows from the comparison principle, Theorem 10.1, applied to the domain
$\mathcal{N}\cap\Omega$, that

\[
w^-(x)\leq u(x)\leq w^+(x)\qquad\text{for all }x\in\mathcal{N}\cap\Omega,
\]
